% ----------------------------------------------------------------------------
\chapter{The lower-dimensional area formula}\label{ch:area}
\module{NoCompromise.Area.AlmostLinear,
        NoCompromise.Area.Cofactor,
        NoCompromise.Area.Formula,
        NoCompromise.Area.GoodPieces,
        NoCompromise.Area.GoodPiecesCover,
        NoCompromise.Area.Graph,
        NoCompromise.Area.Linear,
        NoCompromise.Area.PieceFormula,
        NoCompromise.Area.RankDeficient,
        NoCompromise.Area.Rectifiable,
        NoCompromise.Area.Sphere}

The base supplies the equidimensional change of variables only
(Formalisation note~\ref{fnote:false-friends}).  This chapter derives the
$2$-dimensional area formula in $\R^m$ and its consequence for rectifiable
subsets of $\R^3$.  The decomposition into four elementary lemmas is
deliberate: each is separately formalisable and only the last needs
Rademacher.

\section{Linear distortion}\label{sec:area-linear}

\begin{definition}[$2$-Jacobian]\label{def:jacobian}
For a linear $L:\R^2\to\R^m$ put
$J_2(L):=\sqrt{\det(L^{\mathsf T}L)}$.  For $f:\R^2\to\R^m$ Lipschitz set
$J_2f(x):=J_2(Df(x))$ at differentiability points and $J_2f(x):=0$
elsewhere.
\end{definition}

\begin{lemma}[Linear image measure]\label{lem:linear-distortion}
\uses{def:jacobian}
Let $L:\R^2\to\R^m$ have rank two with singular values $s_1\ge s_2>0$, and
let $B\subset\R^2$ be Borel.  Then
\begin{equation}\label{eq:linear-distortion}
  \Hs^2(LB)=s_1s_2|B|=J_2(L)\,|B| .
\end{equation}
\end{lemma}

\begin{proof}
Choose orthonormal singular bases.  Isometries preserve Hausdorff measure;
the remaining diagonal map is handled first on rectangles and then on Borel
sets by outer regularity.
\end{proof}

\begin{constants}\label{const:almost-linear}
In Lemma~\ref{lem:almost-linear} below, $C_L$ may be chosen uniformly over all $L$ with $\|L\|\le N$ and
$s_2(L)\ge1/N$; it then depends only on $N$.
\end{constants}

\begin{lemma}[Almost-linear maps]\label{lem:almost-linear}
\uses{lem:linear-distortion}
Let $L:\R^2\to\R^m$ have rank two with singular values
$s_1\ge s_2>0$, let $B\subset\R^2$ be Borel, and let $h:B\to\R^m$ satisfy
\begin{equation}\label{eq:almost-linear-hyp}
  |h(x)-h(y)-L(x-y)|\le\varepsilon|x-y|
  \quad(x,y\in B),\qquad\varepsilon<s_2/2 .
\end{equation}
Then $h$ is injective, its inverse on $h(B)$ is Lipschitz, and
\begin{equation}\label{eq:almost-linear}
  \bigl(J_2(L)-C_L\varepsilon\bigr)|B|
  \le\Hs^2(h(B))
  \le\bigl(J_2(L)+C_L\varepsilon\bigr)|B| .
\end{equation}
\end{lemma}

\begin{proof}\uses{lem:linear-distortion}
Injectivity and the Lipschitz inverse follow from
\eqref{eq:almost-linear-hyp} and $\varepsilon<s_2/2$.  Orthogonal projection
onto $L(\R^2)$ followed by $L^{-1}$ differs from the identity by at most
$\varepsilon/s_2$ in Lipschitz norm.  Apply
$\Hs^2(TC)\le\Lip(T)^2\Hs^2(C)$ to this projection and to $h$, and combine
with \eqref{eq:linear-distortion}.
\end{proof}

\section{The area formula in the plane}\label{sec:area-formula}

\begin{definition}[Uniform differentiability piece]\label{def:good-pieces}
\uses{def:jacobian}
Let $f:\R^2\to\R^m$ be Lipschitz.  A Borel set $G$ is a \emph{uniform
differentiability piece} if there are a compact carrier
$G\subset K\subset[-N,N]^2$, an integer $N\ge1$, and a nondecreasing modulus
$\omega:[0,\infty)\to[0,\infty)$ with $\omega(r)\downarrow0$ as
$r\downarrow0$ such that $f$ is differentiable at every $x\in K$ and, for
all distinct $x,z\in K$,
\begin{equation}\label{eq:good-pieces}
\begin{gathered}
  \|Df(x)\|\le N,\qquad s_2(Df(x))\ge1/N,\\
  |f(z)-f(x)-Df(x)(z-x)|\le\omega(|z-x|)|z-x|,\\
  \|Df(x)-Df(z)\|\le\omega(|x-z|).
\end{gathered}
\end{equation}
Every Borel subset of such a piece is again a uniform differentiability
piece with the same compact carrier, $N$, and $\omega$.
\end{definition}

\begin{lemma}[The uniform pieces exhaust the rank-two set]
\label{lem:good-pieces-exhaust}
\uses{def:good-pieces}
The rank-two differentiability points of a Lipschitz map
$f:\R^2\to\R^m$ are, up to a Lebesgue-null set, a countable disjoint union
of uniform differentiability pieces.
\end{lemma}

\begin{proof}
Partition first by integer bounds for $|x|$, $\|Df(x)\|$, and
$s_2(Df(x))^{-1}$.  Lusin's theorem gives countably many compact subsets on
which $Df$ is uniformly continuous, up to a null set.  At a differentiability
point put
\[
 R_x(r):=\sup_{0<|y-x|\le r}
 \frac{|f(y)-f(x)-Df(x)(y-x)|}{|y-x|}.
\]
The supremum may be taken over a countable dense set, so $R_x(r)$ is
measurable, and differentiability says $R_x(r)\to0$.  Egorov on the compact
Lusin pieces, followed by a countable exhaustion, makes this convergence
uniform.  The maximum of the two resulting uniform moduli gives
\eqref{eq:good-pieces} on each compact set.  Finally replace the countable
cover by successive set differences; each resulting Borel piece retains the
compact set from which it came as a carrier.
\end{proof}

\begin{lemma}[Area formula on one uniform piece]\label{lem:area-one-piece}
\uses{def:good-pieces}
Let $G$ be a uniform differentiability piece, let $A\subset G$ be Borel, and
let $q:A\to[0,\infty]$ be Borel.  Then
\begin{equation}\label{eq:area-one-piece}
  \int_Aq(x)J_2f(x)\,dx
  =\int_{\R^m}\sum_{x\in A\cap f^{-1}(y)}q(x)\,d\Hs^2(y).
\end{equation}
\end{lemma}

\begin{proof}\uses{lem:almost-linear}
Let $K,N,\omega$ be as in Definition~\ref{def:good-pieces}.  Partition
$[-N,N]^2$ into half-open grid squares of side $r$.  If $r$ is small enough
that $4\omega(\sqrt2r)<1/N$, then on the intersection of $G$ with any one
square, \eqref{eq:good-pieces} and the last inequality there show that $f$
differs, in the pairwise Lipschitz sense, from any one base derivative
$Df(x_Q)$ by at most $2\omega(\sqrt2r)$.  Lemma~\ref{lem:almost-linear} in
fact makes $f$ injective on $K\cap\overline Q$ and gives, uniformly for every
Borel subset $C$ of $G\cap Q$,
\[
 \left|\Hs^2(f(C))-\int_CJ_2f(x)\,dx\right|
 \le C_N\omega(\sqrt2r)|C|.
\]
Here the replacement of $J_2(Df(x_Q))$ by $J_2(Df(x))$ uses the last line of
\eqref{eq:good-pieces} and the uniform local Lipschitz continuity of $J_2$ on
the compact family $\{L:\|L\|\le N,\ s_2(L)\ge1/N\}$.
The image $f(C)$ is Borel without a descriptive-set theorem: the continuous
injection $f:K\cap\overline Q\to\R^m$ is a homeomorphism onto its compact
image, and $C$ is a Borel subset of that compact domain.

Sum this estimate over the grid squares.  On the target side the sum of the
image indicators is exactly the number of preimages, because the source
squares are disjoint and $f$ is injective within each square.  The total
error is at most $C_N\omega(\sqrt2r)|G|$, which tends to zero.  This proves
\eqref{eq:area-one-piece} for $q=1$.  Applying the same argument to Borel
subsets of $A$ gives simple weights, and monotone convergence gives every
nonnegative Borel $q$.
\end{proof}

\begin{lemma}[The rank-deficient part is null in the image]
\label{lem:rank-deficient}
\uses{def:jacobian}
For $f:\R^2\to\R^m$ Lipschitz, $\Hs^2\bigl(f(\{J_2f=0\})\bigr)=0$.
\end{lemma}

\begin{proof}\uses{lem:good-pieces-exhaust}
Let $Z$ be the set of differentiability points at which $J_2f=0$.  Partition
$Z$ by bounds $|x|\le R$ and $\|Df(x)\|\le N$.  As in the Egorov step of
Lemma~\ref{lem:good-pieces-exhaust}, discard a null set and split the
remainder into countably many Borel pieces $G$.  With
\[
 R_x(r):=\sup_{0<|y-x|\le r}
 \frac{|f(y)-f(x)-Df(x)(y-x)|}{|y-x|},
\]
these pieces may be chosen so that on each of them
\[
 \sup_{x\in G}R_x(r)=:\omega_G(r)\longrightarrow0.
\]
Fix one such $G$ and cover it by grid squares $Q$ of side $r$, choosing
$x_Q\in G\cap Q$ in each occupied square.  For $y\in G\cap Q$, differentiability
at $x_Q$ puts $f(y)$ within $C\omega_G(2r)r$ of
$f(x_Q)+Df(x_Q)(y-x_Q)$.  Since $J_2f(x_Q)=0$, the linear image of $Q$ is a
line segment of length at most $CNr$ (or a point).  With
$\varepsilon_r:=\max\{\omega_G(2r),r\}\downarrow0$, the set $f(G\cap Q)$
is therefore covered by at most $C_N/\varepsilon_r$ balls of radius
$C\varepsilon_rr$.  The sum of their squared diameters is at most
$C_N\varepsilon_rr^2$.

There are $O_R(r^{-2})$ occupied squares.  Hence the $2$-dimensional
Hausdorff content of $f(G)$ at scale $C\varepsilon_rr$ is at most
$C_{N,R}\varepsilon_r\to0$, proving $\Hs^2(f(G))=0$.  The discarded
nondifferentiability and Egorov sets are Lebesgue-null, and their images are
$\Hs^2$-null by the elementary Lipschitz inequality
$\Hs^2(f(B))\le\Lip(f)^2\Hs^2(B)$.  Countable union over all pieces, $N$, and
$R$ proves the claim, including the points where $J_2f$ was defined to be
zero because $f$ is not differentiable.
\end{proof}

\begin{theorem}[Two-dimensional area formula with multiplicity]
\label{thm:area-formula-2d}
\uses{def:jacobian}
Let $A\subset\R^2$ be Borel, $f:\R^2\to\R^m$ Lipschitz with $m\ge2$, and
$q:\R^2\to[0,\infty]$ Borel.  Then
\begin{equation}\label{eq:area-formula-2d}
  \int_Aq(x)\,J_2f(x)\,dx
  =\int_{\R^m}\ \sum_{x\in A\cap f^{-1}(y)}q(x)\ d\Hs^2(y).
\end{equation}
\end{theorem}

\begin{proof}
\uses{lem:area-one-piece,lem:rank-deficient,lem:good-pieces-exhaust}
Lemma~\ref{lem:good-pieces-exhaust} partitions the rank-two differentiability
set, modulo a null set, into countably many disjoint uniform pieces $G_j$.
Apply Lemma~\ref{lem:area-one-piece} to $A\cap G_j$ and sum; on the target
side the sums over $j$ combine into the full preimage sum.  Everything left
in the source is either the null exceptional set discarded in
Lemma~\ref{lem:good-pieces-exhaust}, whose image is $\Hs^2$-null by the
Lipschitz inequality, or lies in $\{J_2f=0\}$, whose image is $\Hs^2$-null by
Lemma~\ref{lem:rank-deficient}.  Both sides therefore ignore it, proving
\eqref{eq:area-formula-2d} with the stated nonnegative Borel weight.
\end{proof}

\section{Rectifiable sets in $\R^3$}

\begin{definition}[Countably $\Hs^2$-rectifiable]\label{def:rectifiable}
$S\subset\R^3$ is \emph{countably $\Hs^2$-rectifiable} if it is
$\Hs^2$-measurable and, up to an $\Hs^2$-null set, contained in a countable
union of images of Lipschitz maps $\R^2\to\R^3$.  Its approximate tangent
plane $T_xS$ is taken in the \emph{chartwise} sense: fix, once for $S$, a
countable family of disjoint Borel pieces $A_j\subset\R^2$ and injective
Lipschitz maps $\psi_j$ with $D\psi_j$ of rank two on $A_j$ whose images cover
$S$ up to an $\Hs^2$-null set (such a family exists by Rademacher's theorem),
and set $T_xS:=\operatorname{range}D\psi_j(z)$ at $x=\psi_j(z)$, $z\in A_j$.
Statements about $T_xS$ refer to this fixed family.
\end{definition}

\begin{theorem}[Area formula on a rectifiable set]\label{thm:area-formula-rect}
\uses{def:rectifiable,def:jacobian}
Let $S\subset\R^3$ be countably $\Hs^2$-rectifiable, $\Phi:\R^3\to\R^3$
Lipschitz, and $q:S\to[0,\infty]$ Borel.  Then at $\Hs^2$-a.e.\ $x\in S$
there is a tangential approximate differential
$D^S\Phi(x):T_xS\to\R^3$, and
\begin{equation}\label{eq:area-formula-rect}
  \int_Sq(x)\,J_2\bigl(D^S\Phi(x)\bigr)\,d\Hs^2(x)
  =\int_{\R^3}\ \sum_{x\in S\cap\Phi^{-1}(y)}q(x)\ d\Hs^2(y).
\end{equation}
If $\Phi$ is differentiable at $x$ (in particular, if it is $C^1$ near
$S$), then $D^S\Phi(x)=D\Phi(x)|_{T_xS}$.
\end{theorem}

\begin{proof}\uses{thm:area-formula-2d,lem:good-pieces-exhaust,
                    lem:rank-deficient}
Refine a countable Lipschitz parametrisation of $S$ into disjoint injective
pieces $\psi_j:A_j\to S$, discarding the parameter sets where
$J_2\psi_j=0$ (null in the image by Lemma~\ref{lem:rank-deficient}).  Apply
Rademacher in the parameter plane to $\psi_j$ and $\Phi\circ\psi_j$.  At a
rank-two point define the unique map $D^S\Phi$ by
\[
 D^S\Phi(\psi_j(z))\circ D\psi_j(z)=D(\Phi\circ\psi_j)(z).
\]
This is independent of the injective chart: on an overlap, both candidates
give the approximate linearisation of $\Phi$ along the common approximate
tangent plane, and uniqueness follows by testing on two independent tangent
density directions.  Apply \eqref{eq:area-formula-2d} to $\psi_j$ and to
$\Phi\circ\psi_j$; the a.e.\ chain rule and cancellation of the common factor
$J_2\psi_j$ give \eqref{eq:area-formula-rect}.  Ordinary differentiability of
$\Phi$ makes the displayed defining identity the usual chain rule and gives
the final assertion.
\end{proof}

\begin{lemma}[Cofactor formula]\label{lem:cofactor}
\uses{def:jacobian}
For linear $L:\R^3\to\R^3$ and a unit vector $\nu$,
\begin{equation}\label{eq:cofactor}
  J_2\bigl(L|_{\nu^\perp}\bigr)=|\cof(L)\,\nu| .
\end{equation}
\end{lemma}

\begin{proof}
Choose an oriented orthonormal basis $a,b$ of $\nu^\perp$ and use
$La\times Lb=\cof(L)(a\times b)$ together with $a\times b=\nu$.
\end{proof}

\begin{corollary}[Graph area]\label{cor:graph-area}
\uses{thm:area-formula-2d,lem:cofactor}
For Lipschitz $f:\R^2\to\R$ and Borel $G\subset\R^2$,
\[
  \Hs^2\bigl(\{(x',f(x')):x'\in G\}\bigr)
  =\int_G\sqrt{1+|\nabla f|^2}\,dx' ,
\]
and for the upward unit normal
$\nu=(-\nabla f,1)/\sqrt{1+|\nabla f|^2}$ one has
$d\Hs^2=\sqrt{1+|\nabla f|^2}\,dx'$ on the graph.
\end{corollary}

\begin{proof}\uses{thm:area-formula-2d,def:jacobian}
Apply \eqref{eq:area-formula-2d} to $x'\mapsto(x',f(x'))$, whose
$2$-Jacobian is $\sqrt{1+|\nabla f|^2}$ by Definition~\ref{def:jacobian}.
\end{proof}

\begin{lemma}[Smooth surfaces are rectifiable]\label{lem:smooth-surface-rectifiable}
\uses{def:rectifiable}
Let $\Sigma\subset\R^3$ be a compact smooth embedded surface.  Then $\Sigma$ is
countably $\Hs^2$-rectifiable, $\Hs^2(\Sigma)<\infty$, and for every
$x\in\Sigma$ the approximate tangent plane of
Definition~\ref{def:rectifiable} coincides with the classical tangent plane
$T_x\Sigma$.
\end{lemma}

\begin{proof}\uses{def:rectifiable,cor:graph-area}
Cover $\Sigma$ by finitely many relatively compact graph subcharts; each is
the image of a smooth map whose derivative is bounded on the closure, hence
of a Lipschitz map from a bounded planar domain.  Then
Corollary~\ref{cor:graph-area} gives its finite area and identifies the
tangent plane.
\end{proof}

\begin{lemma}[Area of a sphere]\label{lem:sphere-area}
\uses{cor:graph-area}
For every $R>0$, $\Hs^2(\partial B_R)=4\pi R^2$; in particular
$\Hs^2(\Sph^2)=4\pi$.
\end{lemma}

\begin{proof}\uses{cor:graph-area}
Write $\partial B_R$ as the union of the two graphs
$x_3=\pm\sqrt{R^2-|x'|^2}$ over $B_R^2$, whose common boundary circle is
$\Hs^2$-null.  These graph functions are not Lipschitz up to
$\partial B_R^2$, so first apply Corollary~\ref{cor:graph-area} on
$B_{(1-1/j)R}^2$, $j\ge2$, and then let $j\to\infty$ by monotone
convergence.  This gives
\[
  \Hs^2(\partial B_R)
  =2\int_{B_R^2}\frac{R}{\sqrt{R^2-|x'|^2}}\,dx'
  =2\cdot2\pi\int_0^R\frac{Rr\,dr}{\sqrt{R^2-r^2}}
  =4\pi R^2 .
\]
\end{proof}

\begin{fnote}[A small gap worth closing early]\label{fnote:sphere-area}
Lemma~\ref{lem:sphere-area} is used by
Lemma~\ref{lem:ball-perimeter} and by
Theorem~\ref{thm:total-curvature-bound}, and it is the kind of fact that is
easy to assume without proof.  Note that it must \emph{not} be derived from
$\Per(B_R)=4\pi R^2$, which is how it would naturally be quoted: that identity
is proved \emph{from} this lemma via
Corollary~\ref{cor:smooth-perimeter}.
\end{fnote}

